\chapter{概率与统计：不确定性也能被精确研究}

\section{一场生日聚会上的赌局}

新学期开学，班里转来了几位新同学，全班一共三十个人。课间有人起哄玩一个游戏：每个人报出自己的生日（只看月和日，不看年份），看看有没有两个人同一天。班长小林信心十足地说：“一年有三百六十五天，咱们才三十个人，撞上的机会太小了。我敢打赌，肯定没有！”同桌小雨却笑了：“我反而敢打赌，很可能有。赌一瓶汽水，来不来？”

同学们纷纷开始报生日。报到第十几个人的时候，教室里忽然爆出笑声：三月十四号出现了两次。再往后报，十一月二号也撞了车。小林目瞪口呆地掏出一瓶汽水的钱，嘴里嘀咕：“这纯粹是运气吧？”

这真的只是运气吗？三十个人里有两个人生日相同，到底是小概率的巧合，还是本该如此？凭直觉，几乎所有人都会站在小林一边。三百六十五对三十，差距悬殊。可是这一章要告诉你：在这类问题面前，人类的直觉几乎总是会输，而数学会赢。我们不仅能把“运气”算成一个精确的数，还能把这个数算给你看。

\section{直觉给出的答案错得离谱}

先别急着看答案，凭直觉猜一猜：一个三十人的班级里，至少有两人生日相同的概率大概是多少？百分之一？百分之十？多数人猜的数字都不到百分之二十。

现在来算一算。直接算“至少两人相同”很麻烦，因为可能是两个人撞、三个人撞、好几对分别撞……情况太杂。数学家遇到“至少”两个字，常用的招数是反过来算：先算“所有人的生日全都不同”的概率，再用 $1$ 减去它。

假设一年有 $365$ 天，每个人的生日等可能地落在任意一天（先不管闰年）。第一个人随便占一天；第二个人的生日要和他不同，有 $\dfrac{364}{365}$ 的机会；第三个人要和前两个人都不同，有 $\dfrac{363}{365}$ 的机会……依次类推，三十个人生日全不相同的概率是
\[
\frac{365}{365}\times\frac{364}{365}\times\frac{363}{365}\times\cdots\times\frac{336}{365}.
\]
这串乘积算出来大约等于 $0.294$。也就是说，“至少两人生日相同”的概率是
\[
1-0.294=0.706,
\]
超过百分之七十！二十三个人时这个概率就超过了一半，五十个人时高达百分之九十七左右。小林的赌局，从一开始就注定大概率要输。

为什么直觉错得这么厉害？因为我们下意识地只拿“三十”和“三百六十五”做比较，却忘了题目比较的不是人数和天数，而是“人对”和天数。三十个人里能配出多少对人呢？第一个人可以和另外二十九人配对，第二个人还可以和后面二十八人配……一共是
\[
\frac{30\times 29}{2}=435
\]
对。四百三十五对，每一对都有撞生日的可能，这就完全不是“三十对三百六十五”的观感了。直觉盯着单个的人，数学盯着人和人之间的关系——这就是它赢的原因。

这个现象太有名了，人们给它起了个名字：生日悖论。注意，它不是真正的逻辑矛盾，而是一条“结论与直觉严重冲突”的定理。数学里凡是叫“悖论”的东西，往往是在提醒我们：直觉需要升级了。

\section{给“可能性”一个数}

刚才的计算里，我们一直在说“概率”，可是概率到底是什么？这个问题曾经困扰了数学家几百年。

回到更简单的场面：掷一枚均匀的硬币，正面朝上的可能性是多少？你会脱口而出：二分之一。为什么？因为结果只有两种，对称，谁也没有理由更偏向谁。一枚均匀的骰子掷出六点呢？六种等可能的结果里占一种，所以是 $\dfrac{1}{6}$。这种想法叫“古典概型”：如果所有基本结果等可能，那么
\[
\text{某事件的概率}=\frac{\text{该事件包含的结果数}}{\text{所有结果的总数}}.
\]
生日问题的计算用的就是这个公式，只不过“结果数”多得要用乘法原理来数。

可是生活中大量的事情没法这样数。明天下雨的可能性是多少？某位球员罚点球命中的可能性是多少？这里既没有有限个等可能的结果，也不能把明天重复一遍。于是我们换一个思路：看频率。让同样的试验重复很多很多次，数一数目标事件出现了几回，用
\[
\text{频率}=\frac{\text{出现的次数}}{\text{试验的总次数}}
\]
去逼近它。掷硬币十次，可能正面出现七次，频率是 $0.7$；掷一千次，频率通常会落在 $0.5$ 附近；掷十万次，几乎一定非常接近 $0.5$。经验告诉我们：随着试验次数增加，频率会稳定在某个数附近，这个稳定的数，就是我们想找的“概率”。

\begin{dingyi}
一个随机试验中，事件 $A$ 的\textbf{概率}，记作 $P(A)$，是这样一个数：当试验在相同条件下独立地重复非常多次时，$A$ 出现的频率会稳定在 $P(A)$ 附近。概率满足三条基本规则：对任何事件 $A$ 有 $0\le P(A)\le 1$；必然发生的事件概率为 $1$；两个不可能同时发生的事件 $A$、$B$，满足 $P(A\text{ 或 }B)=P(A)+P(B)$。
\end{dingyi}

\begin{shuyu}{频率与概率}
直观解释：频率是“已经发生了多少回”，概率是“长远来看会稳定在哪儿”。正式定义：频率是实验中数出来的比值，随每次实验波动；概率是频率在试验次数趋于无穷时逼近的那个数。一个例子：掷均匀硬币一万次，正面出现 $5031$ 次，频率是 $0.5031$，而概率是 $0.5$。
\end{shuyu}

第三条规则有个常用推论：$P(A\text{ 不发生})=1-P(A)$。生日问题里我们正是靠它“反着算”才绕开了繁琐的分类。别小看这一步，从“正面硬算”换成“算反面”，是概率论里最省力的思维方式之一。

还要强调“独立地重复”这几个字。每次掷硬币，硬币不记得上一次的结果。这句话听起来平平无奇，后面我们会看到，不理解它的人会在赌场里付出真金白银的代价。

\section{已知一半信息之后}

现在把问题拧得更紧一点。袋子里有 $3$ 个红球、$2$ 个白球，摸出第一个球不放回，再摸第二个。第二个球是红球的概率是多少？这取决于第一个球摸到了什么——信息变了，可能性就跟着变。这种“在某事已经发生的条件下”的概率，叫条件概率，记作 $P(A\mid B)$，读作“在 $B$ 发生的条件下 $A$ 的概率”。

\begin{dingyi}
设 $P(B)>0$，定义
\[
P(A\mid B)=\frac{P(A\text{ 且 }B)}{P(B)}.
\]
\end{dingyi}

为什么要这样定义？把试验重复 $N$ 次（$N$ 很大），其中 $B$ 大约出现 $N\,P(B)$ 次；在这 $N\,P(B)$ 次里，$A$ 也跟着发生的，大约就是“$A$ 且 $B$”出现的次数 $N\,P(A\text{ 且 }B)$。所以在“$B$ 已经发生”的世界里，$A$ 出现的频率是两者的比值——定义只是把这句大白话写成了符号。

如果 $B$ 的发生丝毫不改变 $A$ 的可能性，即 $P(A\mid B)=P(A)$，就称 $A$ 与 $B$ 相互独立。代入定义式，独立等价于
\[
P(A\text{ 且 }B)=P(A)\,P(B).
\]
先后掷两次硬币，第一次的结果不影响第二次，所以“第一次正面”和“第二次正面”独立，同出正面的概率是 $\dfrac{1}{2}\times\dfrac{1}{2}=\dfrac{1}{4}$。而摸球不放回的情形里，两次摸球并不独立。

条件概率最惊人的应用，是下面这个定理。

\begin{dingli}[贝叶斯公式]
设事件 $B_1,B_2,\dots,B_n$ 把全部可能结果不重不漏地分成几类，且每个 $P(B_i)>0$。则对任一事件 $A$，
\[
P(B_i\mid A)=\frac{P(A\mid B_i)\,P(B_i)}{P(A\mid B_1)P(B_1)+P(A\mid B_2)P(B_2)+\cdots+P(A\mid B_n)P(B_n)}.
\]
\end{dingli}

\begin{proof}
分两步。第一步，分母是什么？事件 $A$ 发生时，必然是和某一个 $B_j$ 一起发生的（因为 $B_j$ 们覆盖了一切可能），而且这些“$A$ 且 $B_j$”两两不能同时出现，于是
\[
P(A)=P(A\text{ 且 }B_1)+P(A\text{ 且 }B_2)+\cdots+P(A\text{ 且 }B_n).
\]
再由条件概率的定义，$P(A\text{ 且 }B_j)=P(A\mid B_j)P(B_j)$，代进去就得到分母那一长串。第二步，看分子。同样由定义，
\[
P(A\text{ 且 }B_i)=P(A\mid B_i)P(B_i)=P(B_i\mid A)P(A).
\]
把最后一个等号两边除以 $P(A)$，再把 $P(A)$ 换成第一步展开的长串，公式成立。
\end{proof}

公式看着吓人，用一个真实味道的例子就立刻清楚了。假设某种疾病在人群中的患病率是 $0.1\%$，即千分之一。有一种检测方法很灵敏：患者检测呈阳性的概率是 $99\%$，健康人检测呈阴性的概率也是 $99\%$（也就是说健康人有 $1\%$ 被误报为阳性）。现在某人检测呈阳性，他真的患病的概率是多少？

直觉会喊：“百分之九十九吧，检测这么准！”算一算。想象一座一百万人的城市：患者约 $1000$ 人，其中约 $990$ 人检测阳性；健康人约 $999000$ 人，其中约 $1\%$、也就是 $9990$ 人被误报阳性。阳性报告一共约 $990+9990=10980$ 份，其中真患者只有 $990$ 人。所以
\[
P(\text{患病}\mid\text{阳性})=\frac{990}{10980}\approx 0.09,
\]
只有百分之九左右！用贝叶斯公式写就是
\[
\frac{0.99\times 0.001}{0.99\times 0.001+0.01\times 0.999}\approx 0.09.
\]

\begin{center}
\begin{tikzpicture}[scale=1.0]
  \node (start) at (0,0) {一百万人};
  \node (ill) at (-3,-2) {患病 $1000$ 人};
  \node (well) at (3,-2) {健康 $999000$ 人};
  \node (illpos) at (-4.5,-4) {阳性 $990$};
  \node (illneg) at (-1.5,-4) {阴性 $10$};
  \node (wellpos) at (1.5,-4) {误报阳性 $9990$};
  \node (wellneg) at (4.5,-4) {阴性 $989010$};
  \draw (start) -- (ill);
  \draw (start) -- (well);
  \draw (ill) -- (illpos);
  \draw (ill) -- (illneg);
  \draw (well) -- (wellpos);
  \draw (well) -- (wellneg);
\end{tikzpicture}
\end{center}

这张图揭示了全部秘密：误报的比例虽小，但健康人的基数太大，假阳性的人数反而十倍于真患者。直觉只盯着“检测准确率 $99\%$”，忘记了患病率这个“先验信息”。贝叶斯公式的全部精神就是一句话：\textbf{新证据要和旧背景一起算}。先验概率 $P(B_i)$ 是旧背景，检测结果 $A$ 是新证据，$P(B_i\mid A)$ 是更新之后的判断。这也正是现代医学筛查、垃圾邮件过滤、搜索引擎拼写纠错背后的共同思想。

\section{平均值能告诉我们什么}

概率回答“会不会发生”，接下来的问题是“平均下来会怎样”。设一个抽奖箱里有 $100$ 张券：$90$ 张“谢谢参与”（奖金 $0$ 元），$9$ 张奖金 $10$ 元，$1$ 张奖金 $100$ 元。抽一张，平均能拿到多少钱？

“平均”在这里的意思不是某一次的结果——单次你只能拿到 $0$、$10$ 或 $100$ 中的一个。它的意思是：如果你抽一万次（每次抽完放回），总奖金大约是多少，分摊到每次是多少。按概率算，一万次里大约有 $9000$ 次 $0$ 元、$900$ 次 $10$ 元、$100$ 次 $100$ 元，总奖金约 $9000\times 0+900\times 10+100\times 100=19000$ 元，平均每次约 $1.9$ 元。写成公式就是
\[
E=0\times 0.9+10\times 0.09+100\times 0.01=1.9\ (\text{元}).
\]

\begin{dingyi}
设随机量 $X$ 能取值 $x_1,x_2,\dots,x_n$，取每个值的概率分别是 $p_1,p_2,\dots,p_n$，则称
\[
E(X)=x_1p_1+x_2p_2+\cdots+x_np_n
\]
为 $X$ 的\textbf{数学期望}（也叫均值）。
\end{dingyi}

期望是“用概率加权的平均”。它有一个立刻能用的推论：如果每张券卖 $3$ 元，而期望奖金只有 $1.9$ 元，那么长期玩下去，玩家平均每局净亏 $1.1$ 元。所有的彩票、赌场游戏都建立在同一笔账上——庄家的期望收益永远为正，玩家的永远为负。期望不能预言下一局的输赢，但它精确地预言了长期的结果。

只看平均值还不够。两个班的平均分都是 $80$ 分，一个班人人都在 $78$ 到 $82$ 之间，另一个班一半是满分、一半不及格，这两个班完全不同。我们需要一个刻画“波动有多大”的数：先算每个取值离期望有多远，把距离平方（平方是为了不让正负抵消），再按概率加权平均，得到
\[
\operatorname{Var}(X)=(x_1-E)^2p_1+(x_2-E)^2p_2+\cdots+(x_n-E)^2p_n,
\]
这叫\textbf{方差}。方差大，说明结果散落得开，随机性强；方差小，说明结果总挤在期望附近。上面那个抽奖的方差很大（大多数人空手而归，少数人拿大奖），而掷骰子点数的方差就小得多。投资里所谓“高风险高收益”，翻译成数学语言就是：期望大，方差也大。

\begin{fanli}
“这个抽奖的期望是 $1.9$ 元，所以我抽一张券，大概能拿到接近两块钱。”——错。单抽一张，你只可能拿到 $0$、$10$ 或 $100$ 元，永远拿不到 $1.9$ 元。期望是关于“长期平均”的陈述，不是关于“下一次”的陈述。把期望当成对单次结果的预言，是最常见的误读之一。
\end{fanli}

\section{掷得足够多，混乱就变成了秩序}

到目前为止，我们一直在用“频率会稳定在概率附近”这句话。它到底是一条铁律，还是只是一种常见的经验？数学家把它变成了定理。

\begin{dingli}[大数定律（直观版）]
把一个随机试验在相同条件下独立地重复 $n$ 次，设每次结果的期望是 $E$。那么当 $n$ 越来越大时，这 $n$ 次结果的算术平均值会以几乎必然的方式逼近 $E$；试验次数越多，平均值偏离 $E$ 超过任何给定小数的概率就越接近 $0$。
\end{dingli}

完整的证明需要更多工具，但它的思想可以摸得着：单次试验的上下波动，在取平均时会互相抵消——有人今天运气极好，就有人运气极差，拉平之后剩下的就是期望。大数定律是保险业的根基：没人知道某位客户明年会不会出车祸，但保险公司面对一百万客户时，总赔付几乎是一个确定的数，于是保费可以精确地定出来。赌场也一样：它不怕某个赌客今晚赢钱，它只怕没有足够多的赌客。大数定律站在它那边。

还有一个更惊人的事实。掷一枚骰子，点数的分布是“平”的：$1$ 到 $6$ 各六分之一，画出来是六根一样高的柱子，毫无钟形可言。可是改掷两枚骰子看点数之和，和为 $7$ 有 $6$ 种凑法，和为 $2$、$12$ 各只有 $1$ 种，分布中间高两头低。掷的骰子越多，和的分布就越接近一条光滑的、对称的钟形曲线。这不是骰子的特殊性质：

\begin{dingli}[中心极限定理（直观版）]
大量相互独立的微小随机因素叠加在一起，不管每个因素本身服从什么分布，它们的和（或平均）的分布都会接近同一条钟形曲线——正态分布。因素越多，逼近得越好。
\end{dingli}

\begin{center}
\begin{tikzpicture}[scale=1.0]
  \draw[->] (-5,0) -- (5,0);
  \draw[->] (0,0) -- (0,3.2);
  \draw[smooth] plot coordinates {(-4.5,0.02) (-3.5,0.08) (-2.5,0.35) (-1.5,1.1) (-0.7,2.2) (0,2.8) (0.7,2.2) (1.5,1.1) (2.5,0.35) (3.5,0.08) (4.5,0.02)};
  \node at (0,-0.4) {$E$};
  \node at (4.6,-0.4) {平均值};
  \node at (0.9,3.0) {次数};
\end{tikzpicture}
\end{center}

这就是为什么身高、测量误差、考试分数、零件尺寸……全都长成差不多形状的钟形：它们每一个都是无数小因素叠加的产物。中心极限定理被称为概率论的“王冠”，因为它解释了为什么混乱的源头各不相同，秩序的形状却惊人地统一。它的证明超出本书范围，难点在于要精确控制“越来越像”的速度，这需要第 12 章介绍的分析工具。

说完秩序，必须回头清算一个顽固的直觉错误。赌场轮盘连续开出了七次红色，围观的人纷纷押黑色：“都红这么多次了，该黑了！”错。轮盘没有记忆，下一次开出红与黑的概率和之前一模一样。这种想法叫\textbf{赌徒谬误}：把“长期会拉平”误当成“短期会补偿”。大数定律只说十万次之后红与黑的比例会接近，从不说接下来十次会往哪个方向“还债”。一九一三年，蒙特卡洛的一家赌场曾连续开出二十六次黑色，押红色“该回了”的赌客们输掉了成堆的法郎——他们每个人都是赌徒谬误的受害者。

\begin{fanli}
“抛硬币连出五次正面，下一次出反面的概率更大。”——错。各次抛掷相互独立，第六次出反面的概率仍然是 $\dfrac{1}{2}$。真正稀有的是“连续六次正面”这整个序列，而不是序列里的某一步。混淆“整段历史”的概率和“下一步”的概率，是赌徒谬误的机关所在。
\end{fanli}

\begin{lianjie}
你在第 9 章学过的极限，在这里换了一副面孔出现：大数定律说的正是“平均值这个数列收敛到期望”，只不过收敛的方式要允许随机性。而中心极限定理里的钟形曲线，它的方程是 $y=\dfrac{1}{\sqrt{2\pi}}e^{-x^2/2}$，里面同时住着 $\pi$ 和 $e$——一个来自圆，一个来自增长与复利。圆、增长、随机性，看似毫不相干的三件事，在这一个公式里会师了。数学各分支之间的暗道，比你想象的多。
\end{lianjie}

\section{用乱撒的点子算出圆周率}

前面所有例子都是“知道概率，预测频率”。现在反过来：利用频率，去算一个本来和随机毫不相干的确定常数。这个思路叫蒙特卡洛方法，名字来自那座以赌场闻名的城市——因为它的燃料就是随机数。

拿一张边长为 $2$ 的正方形纸，在里面画一个半径为 $1$、与四边相切的圆。正方形面积是 $4$，圆面积是 $\pi$。现在闭上眼睛，朝正方形里随机撒点——比如把一把芝麻抛上去，或者在方格纸上蒙眼戳笔尖。假设撒了 $N$ 个点，其中 $M$ 个落在圆内。只要每个点落在正方形里任何位置的机会均等，那么
\[
\frac{M}{N}\approx \frac{\text{圆面积}}{\text{正方形面积}}=\frac{\pi}{4},
\qquad\text{于是}\qquad
\pi\approx \frac{4M}{N}.
\]

\begin{center}
\begin{tikzpicture}[scale=1.6]
  \draw (0,0) rectangle (4,4);
  \draw (2,2) circle (2);
  \fill (0.5,0.8) circle (0.04);
  \fill (1.2,3.4) circle (0.04);
  \fill (2.6,2.9) circle (0.04);
  \fill (3.5,1.1) circle (0.04);
  \fill (1.8,1.5) circle (0.04);
  \fill (2.9,0.6) circle (0.04);
  \fill (0.7,2.4) circle (0.04);
  \fill (3.2,3.6) circle (0.04);
  \fill (2.2,2.1) circle (0.04);
  \fill (1.5,0.4) circle (0.04);
  \fill (3.7,2.5) circle (0.04);
  \fill (0.3,3.7) circle (0.04);
  \node at (2,-0.35) {圆内点数与总点数之比 $\approx \pi/4$};
\end{tikzpicture}
\end{center}

这个办法粗得可爱：撒几百个点，可能得到 $\pi\approx 3.1$；撒几十万个点（当然得请计算机帮忙撒），就能逼近到小数点后两三位。它的精度随点数增长得很慢——误差大致和 $\sqrt{N}$ 成反比，想把精度提高十倍，点数得增加一百倍。所以估算 $\pi$ 它不是最快的办法，它的真正威力在于另一类问题：计算形状极其古怪、根本写不出面积公式的区域的面积，或者高维空间里的体积。物理学里计算粒子穿过材料的轨迹、金融学里评估复杂产品的价格、计算机图形学里渲染光影，背后都是同一招：\textbf{用随机试验把算不出的量“摸”出来}。第 11 章的积分在这里以最意想不到的方式回来了——蒙特卡洛方法本质上就是“用掷骰子做积分”。

还有一个更古典的玩法，叫蒲丰投针。把许多根长度为 $l$ 的针随机抛到画满平行线的纸上，线间距为 $d$（设 $l\le d$）。数学家蒲丰在十八世纪证明：针与某条线相交的概率恰好是 $\dfrac{2l}{\pi d}$。也就是说，反复投针、统计相交的频率 $f$，就能反推出
\[
\pi\approx \frac{2l}{fd}.
\]
一根针、一张纸、一个无聊的下午，就能让圆周率从随机噪声里浮出来。证明的思路要用到积分：针的位置和角度都是随机的，把所有可能的位置按“机会大小”加权累加起来，相交的条件就化成了一个可以积出来的区域。这个思路的细节属于“继续攀登”的部分。

\begin{shiyan}
\textbf{撒点估 $\pi$（家庭版）。}在方格纸上画一个 $10\times 10$ 的大正方形，以一角为圆心、边长为半径，用圆规画一个四分之一圆。然后合上眼，用铅笔尖朝纸上随机戳，每戳一次睁眼记录点落在四分之一圆内还是圆外（戳到纸外就重来）。戳 $100$ 次，若圆内有 $M$ 个点，就用 $\dfrac{4M}{100}$ 估计 $\pi$。和同学把数据合并起来——把五个人的 $500$ 个点合在一起算，比较合并前后的估计哪个更准。你能从中亲眼看到大数定律在工作：点越多，估计越稳。
\end{shiyan}

\begin{toolbox}
本章新添的工具：
\begin{itemize}
  \item \textbf{反向计算}：求“至少发生一次”的概率，先算“一次都不发生”，再用 $1$ 去减。
  \item \textbf{条件概率与贝叶斯公式}：$P(A\mid B)=\dfrac{P(A\text{ 且 }B)}{P(B)}$；新证据要和先验背景一起算。
  \item \textbf{独立性检验}：$P(A\text{ 且 }B)=P(A)P(B)$ 成立才叫独立，别凭感觉。
  \item \textbf{期望与方差}：$E(X)$ 管长期平均，$\operatorname{Var}(X)$ 管波动大小。
  \item \textbf{大数定律}：次数够多，平均值就逼近期望——它保证长期，不承诺下一局。
  \item \textbf{蒙特卡洛方法}：用随机试验的频率逼近面积、常数和积分。
\end{itemize}
\end{toolbox}

\section{当统计开始骗人}

学会了工具，更要学会识别陷阱。除了生日悖论和赌徒谬误，还有几个专门坑害“半懂不懂”的人的机关。

第一个叫幸存者偏差。二战时，统计学家研究返航轰炸机上的弹孔分布，机翼上弹孔最多。军官们主张加固机翼，而数学家建议加固的恰恰是弹孔最少的发动机部位——因为发动机中弹的飞机根本没飞回来，你看到的全是幸存者。只看得见回来的飞机，数据就先天地缺了一半。今天的“读书无用论”（举几个没上学的富翁当证据）、“成功学课程”（只采访活下来的公司），犯的是同一个错误。

第二个叫平均数幻觉。“某公司员工平均年薪五十万”——可能老板一个人年薪五千万，其他人都在五万上下。期望对极端值极其敏感，这时“中位数”（排最中间的那个数）才是更诚实的描述。下次看到任何“平均”，先问一句：这是哪种平均？被平均的又是谁？

第三个是把相关当因果。冰淇淋销量和溺水人数同步上升，难道是冰淇淋导致溺水？其实背后有个共同原因：夏天到了。两件事一起涨，可能是一个涨导致另一个涨，可能方向相反，也可能纯属第三者作祟甚至巧合。统计能发现关联，但确认因果需要对照实验或更深的机制分析。这是数据时代每个公民的必修课。

\begin{pandeng}
继续向前，这些关键词在等你：\textbf{正态分布}（钟形曲线的精确方程与“$3\sigma$ 法则”）、\textbf{假设检验}（“这个差异是巧合还是真的有效”的定量裁判）、\textbf{大数定律的严格证明}（切比雪夫不等式是入场券）、\textbf{马尔可夫链}（“明天只取决于今天”的随机过程，网页排序的数学引擎）、\textbf{信息论}（用概率度量“惊讶程度”，单位是比特）。竞赛和高中课程里，还会遇到用递推求概率的漂亮问题，比如著名的“赌徒破产问题”。
\end{pandeng}

\section*{思考题}

\textbf{观察题}

1. 找一副扑克牌，洗牌后一张张翻开，记录前两张牌同花色的情况，重复二十次，估计“前两张同花色”的频率。再算一算理论概率（提示：第一张任意，第二张与第一张同花色的概率只取决于还剩几张同花色的牌），与你的频率比较。

2. 收集家里一周的用电量数据（电表读数之差），算平均值，再想想：哪个因素会让某天的读数偏离平均最远？这算不算方差的来源？

\textbf{证明题}

1. 证明：若事件 $A$ 与 $B$ 相互独立，则“$A$ 不发生”与“$B$ 不发生”也相互独立。（提示：用 $P(A\text{ 不发生})=1-P(A)$ 展开 $P(\text{都不发生})=1-P(A)-P(B)+P(A)P(B)$，再做因式分解。）

2. 把生日问题换一种问法：要多少人才能保证“至少两人生日相同”的概率超过 $99\%$？给出你的估算过程。（提示：不必精确求解，逐个试 $40$、$50$、$60$ 人时的“全不同”概率，利用对数或分步乘积估计。）

\textbf{挑战题}

1. 蒙提霍尔问题：三扇门后有一辆汽车和两只山羊。你选定一扇门后，主持人（他知道答案）打开另外两扇中藏有山羊的一扇，然后问你要不要换门。换门赢得汽车的概率是多少？用条件概率或枚举给出论证，并解释为什么“还剩两扇门所以各二分之一”是错的。（提示：关键是主持人不是随机开门——他的动作携带了信息。）

2. 设计一个蒙特卡洛实验来估计 $\sqrt{2}$ 的近似值，只需用到撒点和面积比较。（提示：$\sqrt{2}$ 是某个简单几何图形的边长或对角线长；也可以考虑比较曲线 $y=x^2$ 附近的面积。）

\section*{通往下一章的门}

这一章里，我们好几次在紧要关头说出同一句话：“撒几十万个点”“掷十万次”“重复非常多回”。数学给出了公式，可是谁来做这几十万次枯燥的重复？人手做不到，人脑会出错，而有一台机器天生就爱干这个：它每秒能掷几十亿次骰子，从不喊累，从不抱怨。

蒙特卡洛方法真正改变世界，正是在电子计算机诞生之后——第一台通用电子计算机投入使用时，数学家们最先让它算的，就是用随机数模拟中子在材料里的穿行。可是新的问题立刻来了：计算机产生的“随机数”其实是由确定公式算出来的，它还算真正的随机吗？哪些问题计算机眨眼就能模拟，哪些问题给它一万年也算不完？下一章，我们推开数学与计算机之间的那扇门。
